\documentclass[journal]{IEEEtran}

\usepackage{cite}
\usepackage{amsmath,amssymb,amsfonts}
\usepackage{graphicx}
\usepackage{textcomp}
\usepackage{xcolor}
\usepackage{url}
\usepackage{booktabs}

\begin{document}
	
	\title{Embedded Machine Learning-Based Voltage Fingerprinting for Automotive Cybersecurity: A Survey}
	
	\author{Goki~\IEEEmembership{Student}%
		\thanks{Manuscript in preparation.}}
	
	\markboth{Preprint / Draft for Review}{}
	
	\maketitle
	
	\begin{abstract}
		Placeholder abstract -- to be finalized once all sections are complete.
	\end{abstract}
	
	\begin{IEEEkeywords}
		Automotive cybersecurity, Controller Area Network (CAN), voltage
		fingerprinting, embedded machine learning, intrusion detection,
		electronic control unit (ECU) identification.
	\end{IEEEkeywords}
	
	\section{Introduction}
	
	Modern vehicles have moved from simple mechanical assemblies into
	extensively networked cyber-physical platforms. A typical vehicle now
	integrates between 70 and 100 Electronic Control Units (ECUs) exchanging
	roughly 2,500 internal signals over intra-vehicle networks
	\cite{ref6,ref8,ref9}, and industry sources place the ECU count as high as
	``more than a hundred'' in some modern platforms, with that number expected
	to keep growing \cite{ref10}. These ECUs are grouped into functional
	subnetworks -- powertrain, chassis, body, infotainment -- interconnected
	through gateways using a mix of protocols: the Local Interconnect Network
	(LIN) for low-speed, non-safety-critical functions such as window and seat
	motors; FlexRay for high-bandwidth, deterministic applications like
	steering-angle sensing and safety radar; Automotive Ethernet and the Media
	Oriented Systems Transport (MOST) protocol for infotainment; and the
	Controller Area Network (CAN) bus, which remains the dominant standard for
	powertrain and body electronics owing to its low cost, mature tool-chain,
	and strong fault tolerance \cite{ref4,ref6,ref9,ref10}. CAN itself is a
	multi-master broadcast protocol developed by Robert Bosch GmbH in the
	early 1980s, standardized under ISO 11898-1:2015, and extended in 2012
	with the higher-throughput CAN FD variant (up to 5 Mbps and a 64-byte
	payload, versus 8 bytes for classical CAN) \cite{ref10}. Every CAN frame
	carries an 11- or 28-bit arbitration identifier, a data field of up to 8
	bytes, a 15-bit CRC checksum, and fixed ACK/EOF delimiter bits; nodes
	resolve simultaneous transmission through non-destructive bitwise
	arbitration, in which the frame with the numerically lowest identifier --
	and therefore highest priority -- wins the bus without any collision or
	retransmission delay \cite{ref10}. This design is elegant for real-time
	control, but it was built for a closed, physically isolated network: CAN
	provides no message authentication, no encryption, and no means of
	verifying the true origin of a frame, so once an adversary gains bus
	access -- physically or, increasingly, remotely -- nothing in the
	protocol itself prevents that adversary from injecting or spoofing frames
	that every other node on the bus will treat as legitimate
	\cite{ref4,ref6,ref10}.
	
	Embedded ML-based voltage fingerprinting is the product of a progressive
	convergence of three separate research threads: the maturation of vehicle
	automation and connectivity itself, the parallel rise of
	automotive-specific cybersecurity engineering and regulation, and a long
	lineage of CAN intrusion-detection research that has steadily moved from
	purely cryptographic and message-content countermeasures toward
	physical-layer signal analysis. An in-depth understanding of this
	technical progression is essential to situate where voltage fingerprinting
	fits, and why it has emerged now rather than a decade ago. The following
	section traces this historical trajectory before turning to the standards
	landscape that makes embedded, real-time detection a regulatory necessity
	and the specific research gaps that motivate a physical-layer approach.
	
	\subsection{History}
	
	The trajectory of automotive cybersecurity research is closely tied to
	the trajectory of vehicle automation itself. In 2014, the Society of
	Automotive Engineers (SAE) published the J3016 standard, defining six
	levels of driving automation from Level 0 (no automation) to Level 5
	(full automation); the standard was significantly revised in
	collaboration with ISO in 2021 and is now the industry's most-cited
	reference on autonomous-vehicle capability, endorsed by UNECE
	WP.29/GRVA, the U.S. NHTSA and DOT, and the California DMV as a basis for
	further AV regulation \cite{ref1,ref2}. Commercial deployment followed
	close behind the standardization effort. Google fitted an experimental
	vehicle with lane-keeping and lane-changing capability in 2010, the same
	year Toyota debuted the Lexus LS460L with an advanced parking-assist
	system \cite{ref2}. General Motors followed in 2017 with its IoV-based
	ENV series of self-parking smart cars, and Tesla progressively rolled out
	its Autopilot system -- culminating in Autopilot 8.1 -- before offering
	``full self-driving'' functionality to private owners \cite{ref2}.
	Waymo, sharing a parent company with Google, became the first provider
	to offer public driverless robo-taxi rides, launching in Phoenix,
	Arizona, in 2020 \cite{ref1}. Underpinning much of this progress was the
	parallel maturation of deep learning: Hinton's deep belief network (DBN)
	in the mid-2000s catalyzed a wave of deep-structured learning
	architectures that autonomous-vehicle developers subsequently adopted for
	object detection, image classification, semantic segmentation, and
	traffic-sign recognition, tasks for which abundant high-quality driving
	datasets made deep models particularly attractive relative to classical
	ML \cite{ref2}.
	
	Bibliometric analysis of the Web of Science database shows that the
	first dedicated automotive-cybersecurity article was published in 2012,
	after which output accelerated sharply -- of 801 studies identified under
	the search terms ``cybersecurity''+``vehicle''/``automotive,'' roughly
	two-thirds were published within just the three years preceding the
	survey, and 133 were retained after PRISMA-style screening for relevance
	and citation impact \cite{ref3}. This growth was not merely academic
	curiosity; it tracked a parallel history of real intrusions and
	near-misses. Early attack research showed that adversaries with physical
	access could inject malicious messages directly onto the CAN bus, read
	and modify ECU memory through interfaces such as CD players, USB ports,
	and the OBD-II diagnostic port, and even extract ECU security keys
	\cite{ref4}. As wireless connectivity -- Bluetooth, Wi-Fi, cellular, LTE,
	5G -- became standard, many of these previously physical-access-only
	attacks became remotely exploitable. Researchers demonstrated the
	ability to hack vehicle key fobs and replay radio signals to unlock doors
	without defeating any security key at all \cite{ref4}; hobbyists and
	researchers alike showed that tire-pressure monitoring systems (TPMS)
	could be spoofed to display false readings and trigger bogus warnings
	\cite{ref4}; and a team reportedly compromised a Tesla Model S remotely
	from 12 miles away \cite{ref4}. Researchers at Keen Security Labs
	separately demonstrated a remote exploit of a Tesla Model S through its
	camera system \cite{ref2}. Kaspersky Lab researchers exposed software
	vulnerabilities in EV charging stations (EVCSs) that left them accessible
	to unauthorized hackers \cite{ref1}, and Cai et al. uncovered multiple
	vulnerabilities across several BMW models, including the ability to
	compromise CAN-connected ECUs over a wireless link \cite{ref4}. The most
	consequential single event was arguably Miller and Valasek's 2015 remote
	compromise of a Jeep Cherokee traveling on a highway, which forced a
	recall of 1.4 million vehicles and, per Luo et al., ``brought public
	attention to automotive cybersecurity'' for the first time at scale
	\cite{ref5}. Safety consequences followed directly from the security
	failures: a Tesla operating in Autopilot mode was involved in a fatal
	crash in Florida on May 7, 2016, and an Uber self-driving test vehicle
	struck and killed a pedestrian in Tempe, Arizona, on November 20, 2018 --
	itself the 38th crash involving that Uber AV fleet \cite{ref2}.
	
	Girdhar et al. later observed that existing post-collision investigation
	frameworks, including the California AV collision reporting process
	(which as of February 2022 cited causal factors such as incorrect
	traffic-signal perception, degraded localization, map discrepancy, and
	reduced visibility from occlusion) and ISO 12353-1 for conventional
	accident analysis, catalogue an extensive list of physical and
	perceptual causal factors but contain no provision for cyber-induced
	causes at all -- a gap that motivated their proposed 5Ws-and-1H
	digital-forensic investigation model, built on STRIDE threat modeling and
	a stochastic anomaly-detection layer, specifically to reconstruct crash
	scenarios in which a cyberattack on the vehicle, the charging
	infrastructure, or the cooperative intelligent transport system (C-ITS)
	may have been a contributing cause \cite{ref1}. On the detection side,
	the field's earliest systematic reviews began cataloguing defenses in
	parallel with cataloguing attacks. Al-Jarrah et al.'s 2019 survey of
	intra-vehicle IDS work reviewed 42 detection schemes and organized them
	by detection technique, feature-selection method, evaluation dataset,
	targeted attack type, and benchmark model \cite{ref6}. Abu Elkhail et al.
	separately catalogued the malware-defense landscape into five categories
	-- signature-based, behavior-based, heuristic-based, cloud-based, and
	machine-learning-based detection -- while noting that malware could enter
	a vehicle through wireless roadside networks, Wi-Fi hotspots, cellular
	connectivity, or even infected consumer devices such as phones and
	laptops physically or wirelessly tethered to the car \cite{ref4}. Bozdal
	et al.'s comprehensive CAN vulnerability assessment catalogued dozens of
	experimentally implemented attacks on both passenger cars and
	heavy-duty vehicles, and traced the shift in defensive research from
	purely cryptographic countermeasures toward reactive, traffic-monitoring
	intrusion detection systems (IDS) as a complementary layer of defense
	precisely because CAN's built-in error-handling features -- bit-level
	monitoring, the 21-bit CRC field, and the error-confinement mechanism
	that forces a faulty node ``bus off'' once its transmit-error counter
	exceeds 255 -- were designed for reliability, not security, and do
	nothing to stop a well-formed but malicious frame \cite{ref10}.
	
	\subsection{Rationale}
	
	The rationale for continued, and increasingly urgent, investment in
	automotive cybersecurity is grounded in a widening gap between the pace
	of vehicle automation and the pace of safety validation. Girdhar et al.
	report that autonomous vehicles currently average 9.1 crashes per million
	miles driven, more than double the 4.1 crashes per million miles for
	conventional human-driven vehicles \cite{ref1} -- a statistic that
	undercuts the common assumption that automation is inherently safer, and
	one that places extra weight on eliminating cyber-induced failure modes
	specifically. Adu-Kyere et al. add a market-scale dimension to this
	urgency, citing projections of 115 million connected vehicles by 2025
	alongside reports that the frequency of attacks on modern cars increased
	by 225\% over the preceding three years \cite{ref12} -- figures that
	describe an attack surface growing considerably faster than the fleet
	itself.
	
	Regulatory bodies have responded, but on a lag, and the standards
	landscape itself has evolved in visible stages. Historically, functional
	safety and functional security for conventional road vehicles were
	governed separately by ISO 26262 and SAE J3061, neither of which was
	written with autonomous-specific functionality in mind \cite{ref1}. SAE
	published J3061 in 2016, establishing the now-foundational principle that
	cybersecurity must be designed into a feature from the outset rather than
	appended at the end of development \cite{ref3}. In 2021, the UNECE's
	World Forum for Harmonization of Vehicle Regulations published UN
	Regulation No. 155, requiring every OEM to operate a formal Cyber
	Security Management System (CSMS) addressing organizational processes,
	responsibilities, and governance across the vehicle lifecycle, and to
	demonstrate that risk-management process when seeking vehicle type
	approval \cite{ref3}. ISO/SAE 21434:2021 subsequently codified the
	detailed engineering requirements for managing cybersecurity risk across
	concept, product development, production, operation, maintenance, and
	decommissioning phases of electrical/electronic (E/E) systems, becoming
	mandatory for all newly manufactured vehicles from July 2024
	\cite{ref3}. A recurring theme across the surveyed standards literature
	is that this new Cyber Security Management System must be integrated
	into an automotive software development lifecycle that was not
	originally built to accommodate it -- Kifor and Popescu note that, in
	the initial implementation phase, CSMS processes are typically bolted on
	as an ``add-on'' to existing tools, mirroring how ISO 26262
	functional-safety requirements were only gradually absorbed into
	mainstream development practice over the preceding decade \cite{ref3}.
	Complementary standards round out the landscape: UL4600 and the Safety
	of the Intended Functionality (SOTIF) framework address autonomous-system
	safety cases where no human is available to intervene on failure, while
	ISO 34501--34503 and IAMTS/ASAM guidance extend the SAE J3016 taxonomy
	into concrete testing and simulation practice \cite{ref1,ref2}. Yet even
	with this apparatus in place, Luo et al.'s systematic review of
	automotive cybersecurity \emph{testing} practice found that where
	ISO/SAE 21434 and UNECE WP.29 (R155/R156) explicitly mandate security
	verification and validation, the standards provide no specific technical
	guidance on how that testing should actually be carried out -- the
	guidance is described as ``high-level,'' making ``the actual execution
	of the test difficult'' for OEMs and suppliers who must independently
	decide on cost, automation level, and safety of their testing process
	\cite{ref5}. This is precisely the gap that motivates a stronger
	empirical and technical literature: standards establish that testing and
	detection are required, but not how to build detection mechanisms that
	are simultaneously accurate, real-time, and deployable on
	resource-constrained ECUs.
	
	\subsection{Motivation}
	
	Despite more than a decade of concentrated IDS research, several
	converging observations from the surveyed literature motivate a shift
	toward physical-layer, embedded ML-based detection specifically. First,
	the field remains fragmented at the taxonomic level. Karopoulos et al.'s
	meta-survey compares twelve prior IDS surveys published between 2018 and
	2021 and finds persistent disagreement over even basic categorization:
	some surveys scope themselves to CAN alone while others span LIN,
	FlexRay, Ethernet, and MOST; some provide explicit performance comparison
	across reviewed works while others do not; and the number of underlying
	works reviewed ranges from as few as 5 to as many as 42 across otherwise
	comparable surveys \cite{ref9}. Beneath this inconsistency, however, a
	recurring four-way taxonomy does emerge across the literature:
	message-parameter-based detection (operating at the MAC layer),
	information-theory-based detection (typically exploiting entropy),
	machine-learning-based detection, and fingerprinting-based detection
	operating on the physical layer via clock or transmission-energy
	variation \cite{ref9}. One particularly telling data point emerges
	directly relevant to this review's focus: of the fingerprinting-based
	VIDS approaches surveyed, four out of five achieved greater than 96\%
	accuracy -- the highest reliability band among all IDS categories
	reviewed -- but three of those five required the addition of extra
	dedicated hardware to the vehicle, undercutting their practicality for
	retrofit or low-cost deployment \cite{ref9}. This tension between high
	accuracy and hardware overhead is exactly the design space that
	voltage-fingerprinting on embedded, resource-constrained platforms must
	resolve.
	
	Second, the machine-learning side of the taxonomy has itself grown
	considerably more sophisticated without fully closing the deployment
	gap. Ghadi and Khan Mohd's 2026 deployment-focused survey organizes
	learning-based CAN IDS into ML-based, DL-based, and hybrid categories:
	classical ML methods such as Support Vector Machines, Decision Trees,
	Random Forests, K-Nearest Neighbors, XGBoost, and Naïve Bayes offer model
	transparency and low computational overhead well suited to
	resource-limited ECUs, but their dependence on manually engineered
	features limits generalization to unseen attacks; DL methods -- Recurrent
	Neural Networks, LSTM networks, Convolutional Neural Networks, and
	Transformer-based models -- overcome this by automatically extracting
	feature and sequence-based patterns from raw CAN traffic and achieve
	state-of-the-art accuracy on complex attack vectors, at the cost of
	heavier compute; and hybrid ML/DL approaches attempt to combine the two
	\cite{ref7}. Across studies published between 2019 and 2025 surveyed
	under this taxonomy, the authors found that ``claims related to
	real-time performance or lightweight deployment are frequently made
	without sufficient practical validation,'' and the survey found it
	necessary to introduce its own explicit criteria for what should count
	as evidence of ``lightweight,'' ``real-time,'' or ``deployable'' --
	namely ECU-oriented runtime analysis, embedded evaluation, or
	real-vehicle experimentation -- because so few prior works substantiated
	these claims on genuine embedded hardware \cite{ref7}. Almehdhar et al.
	corroborate this from a different angle, situating CAN threats within a
	three-layer Autonomous Vehicular Sensing Communication Control
	(AutoVSCC) framework: the sensing layer is vulnerable to spoofing and
	eavesdropping on inertial or radar sensors, GPS jamming, TPMS
	interception, and ultrasonic-sensor manipulation; threats here propagate
	upward through the communication layer (eavesdropping, message
	manipulation, unauthorized infotainment/telematics access) and
	ultimately reach the control layer, corrupting real-time functions such
	as automated steering, lane-change maneuvers, and braking \cite{ref8}.
	This layered propagation matters directly for detection-system design:
	an IDS that inspects only message content at the communication layer is
	structurally blind to attacks that originate at the sensing layer and
	only manifest as anomalies in the underlying physical signal.
	
	Third, the threat landscape for AI-enabled autonomous functions
	specifically has opened an entirely separate motivational thread.
	Girdhar, Hong, and Moore's systematic review of adversarial attacks on
	autonomous driving finds that deep-learning models used for object
	detection, image classification, and semantic segmentation are
	vulnerable to deliberately crafted perturbations -- patterns
	imperceptible to human drivers but sufficient to reliably deceive the AI
	model -- and that this vulnerability has so far been demonstrated
	primarily on image-classification models, leaving the effectiveness of
	both attacks and defenses against regression-style autonomous-driving
	models comparatively unexplored \cite{ref2}. Because DL models operate
	largely as black boxes during inference, with limited interpretability
	or explainability, conventional defenses developed for other domains do
	not transfer cleanly, and the authors note that even research using
	explainable-AI methods (such as logic learning machines or support
	vector data description) has been needed to correctly flag adversarial
	attacks that classical algorithms missed entirely \cite{ref2}. Taken
	together, these findings motivate embedded ML-based voltage
	fingerprinting as a response to a well-documented, specific gap:
	existing IDS work already shows that physical-layer signal
	characteristics deliver the highest accuracy of any detection category,
	but has not yet resolved how to deliver that accuracy without extra
	hardware, on real embedded automotive processors, in real time, and
	largely independent of the content-level vulnerabilities that plague
	both traditional message-based IDS and AI perception models alike.
	
	\subsection{Scope}
	
	This survey focuses on embedded machine learning-based voltage
	fingerprinting as a physical-layer complement to conventional,
	message-content-based CAN intrusion detection. It reviews the
	fundamentals of CAN protocol security and the regulatory standards now
	governing it, examines the architectures and taxonomies of ML-based
	fingerprinting approaches, benchmarks their detection performance and
	embedded deployment feasibility against the lightweight/real-time
	evidentiary standard raised by recent deployment-focused surveys, and
	surveys the open challenges -- fabrication/hardware overhead, dataset
	standardization, and cross-vehicle generalization -- and emerging
	directions in the field. It extends the anchor work of Dini
	\textit{et al.} on embedded ML-based voltage fingerprinting for CAN bus
	ECU identification \cite{refDini} into a broader, systematically
	organized review of the surrounding literature spanning historical IDS
	taxonomies,
	adversarial-ML threats to AI-enabled vehicle functions, and the standards
	landscape that now mandates -- without yet fully specifying -- how
	automotive security must be tested and proven.
	
	\section{Fundamentals of Vehicular Network Security and Trust Establishment}
	
	Modern connected vehicles are built on a layered communication fabric
	rather than a single network. The Controller Area Network (CAN) bus
	forms the foundational layer, linking ECUs through broadcast,
	identifier-arbitrated frames, as described in Section I. Automotive
	Ethernet is emerging as a parallel high-bandwidth layer for
	sensor-fusion and autonomy workloads \cite{ref28}. Above both sits the
	outward-facing V2X/Internet-of-Vehicles (IoV) layer, connecting the
	vehicle to roadside infrastructure, edge servers, other vehicles, and
	the cloud \cite{ref13,ref14,ref15,ref16,ref17,ref18,ref19,ref20,ref30}.
	Unlike conventional enterprise IT networks, all three layers interface
	directly with real-time physical control loops, so cybersecurity
	mechanisms that work in IT contexts -- heavy encryption, deep packet
	inspection, centralized certificate revocation -- frequently cannot be
	applied without redesign. This constraint is what has driven the shift
	toward lightweight, hardware-rooted, and decentralized trust mechanisms
	purpose-built for vehicular limits, rather than borrowed wholesale from
	IT security.
	
	Recent work indicates that three closely related factors govern the
	effective operation of automotive cybersecurity architectures. (i) The
	network's protocol topology and bandwidth budget determine how much
	authentication payload, key material, and cryptographic computation a
	link can carry without violating real-time latency constraints
	\cite{ref21}. (ii) The physical and cryptographic trust-establishment
	principle in use -- asymmetric or post-quantum key exchange, physical
	unclonable functions (PUFs), or blockchain-anchored consensus --
	determines how communicating entities verify one another without a
	single point of failure \cite{ref13,ref14,ref15,ref16,ref17,ref18,ref19,ref20,ref22}.
	(iii) The risk-assessment and access-control methodology applied --
	threat analysis and risk assessment (TARA) frameworks, formal
	verification, or selective/targeted authentication policy -- determines
	how residual threats are quantified, prioritized, and mitigated across
	the vehicle's development and operational lifecycle
	\cite{ref29,ref30}.
	
	\begin{figure}[t]
		\centering
		\fbox{\parbox[c][5.5cm][c]{8cm}{\centering\itshape
				Placeholder --- diagrammatic overview of the three-layer vehicular
				network stack (CAN / Automotive Ethernet / cloud-assisted V2X-IoV)
				and their associated trust boundaries, to be inserted.}}
		\caption{Layered vehicular network architecture and associated
			trust-establishment boundaries.}
		\label{fig:layered_arch}
	\end{figure}
	
	\subsection{Network Architectures Relevant to Trust Integration}
	
	The topology of the network directly constrains how much authentication
	and cryptographic overhead a security mechanism can carry without
	violating the link's latency budget. Structural parameters -- frame
	payload size, arbitration scheme, available bandwidth, and the number
	of interconnected ECUs, RSUs, or cloud endpoints -- jointly determine
	bus load, message latency, and exposure to spoofing. Structural
	characterization is therefore a prerequisite for designing any
	deployable defense. Vehicular network architectures can be broadly
	classified by backbone protocol into CAN-based, Automotive-Ethernet-based,
	and cloud-assisted V2X/IoV architectures.
	
	CAN-based architectures carry an 8-byte classical payload per frame and
	no native authentication, so any added security mechanism competes
	directly with control-message bandwidth. Wang et al. \cite{ref21} show
	this trade-off concretely: their dynamic compression and per-ECU
	key-rotation scheme reduces bus load by roughly 8\% on average while
	keeping single-message authentication computation in the
	sub-millisecond range at automotive ECU clock speeds -- evidence that
	CAN security design is fundamentally a bandwidth-allocation problem,
	not only a cryptographic-strength problem. The bus-load overhead
	imposed by an added authentication mechanism can be expressed as
	\begin{equation}
		\rho = \frac{N_{\text{auth}} \times L_{\text{auth}}}{B \times T}
		\label{eq:busload}
	\end{equation}
	where $N_{\text{auth}}$ is the number of authentication frames
	transmitted per window, $L_{\text{auth}}$ is the authentication
	payload length in bits, $B$ is the bus bandwidth, and $T$ is the
	observation window.
	
	Automotive-Ethernet-based architectures relax the payload constraint but
	inherit IT-style network threats. Wang et al. \cite{ref28} characterize
	In-Vehicle Ethernet as the backbone of next-generation in-vehicle
	communication and show that data theft, tampering, and injected traffic
	on protocols such as AVTP require traffic-pattern-based intrusion
	detection rather than per-frame cryptographic authentication, reporting
	a 99.7\% anomaly detection rate with their weighted-histogram method.
	
	Cloud-assisted V2X/IoV architectures extend the trust boundary outward
	again, to RSUs, edge servers, and cloud management centers
	\cite{ref13,ref14,ref15,ref16,ref17,ref18,ref19,ref20}. Sheik et al.
	\cite{ref30} show that this extension changes the threat surface
	qualitatively: complex hardware--software interaction and
	heterogeneous protocol stacks mean that off-the-shelf risk-scoring
	templates poorly capture Cloud-Assisted Connected and Autonomous
	Vehicle (CCAV) threats, motivating their STRIDE/DREAD-based
	reassessment. Risk under DREAD is scored as
	\begin{equation}
		\mathrm{Risk}_{\mathrm{DREAD}} = \frac{D + R + E + A + \mathrm{Disc}}{5}
		\label{eq:dread}
	\end{equation}
	where $D$, $R$, $E$, $A$, and $\mathrm{Disc}$ are the Damage,
	Reproducibility, Exploitability, Affected-users, and Discoverability
	sub-scores \cite{ref30}.
	
	Taken together, architecture, trust principle, and risk quantification
	combine to determine the achievable security posture of a vehicular
	network: none of the three is sufficient alone, and design choices in
	one constrain what is achievable in the others.
	
	\subsection{Mechanisms of Threat and Trust Interaction}
	
	In vehicular networks, entities interact through several distinct
	threat and trust-establishment mechanisms, depending on the network
	layer and the adversary's capability. These include message injection
	and spoofing, eavesdropping and man-in-the-middle interception,
	resource-exhaustion attacks on offloading infrastructure, cryptographic
	key-agreement--based mutual authentication, and decentralized
	ledger-based consensus trust. Each mechanism shapes a different class
	of defense.
	
	Message injection and spoofing is the fundamental threat mechanism on
	broadcast media such as CAN, where any node with bus access can
	transmit frames impersonating another ECU; Wang et al. \cite{ref21}
	address this directly through per-message dynamic keying rather than
	static authentication, since a static key is trivially replayable on a
	broadcast bus.
	
	Eavesdropping and man-in-the-middle interception dominate on V2X
	wireless links, where Dai et al. \cite{ref13} note that classical
	RSA/ECC-based key exchange is vulnerable both to passive interception
	today and, prospectively, to quantum-capable adversaries --
	motivating their hybrid Kyber+Saber post-quantum handshake -- while
	Escapa Gord\'on et al. \cite{ref15} address the same interception risk
	through identity-based cryptography keyed to the Vehicle Identification
	Number (VIN), removing the certificate-exchange step that is itself a
	common interception target.
	
	Resource-exhaustion attacks target the newer V2X computing-offload
	surface: Liu et al. \cite{ref14} show that an attacker impersonating
	either a requesting vehicle or an edge server can exhaust server
	compute resources or exfiltrate offloaded task data, requiring
	four-party authentication (vehicle, RSU, edge server, cloud management
	center) before any task offload is accepted.
	
	Cryptographic key-agreement--based mutual authentication is the active
	countermeasure mechanism common to \cite{ref13,ref14,ref15,ref22}: two
	or more parties jointly derive a shared session key such that neither
	could have derived it alone, closing the injection and interception
	mechanisms above. Madububa Mbachu et al. \cite{ref22} root this process
	in hardware via a Physical Unclonable Function (PUF), deriving the key
	material from FPGA-level manufacturing variation rather than a stored
	secret, which removes key-extraction-by-firmware-dump as an attack
	vector.
	
	Decentralized ledger-based consensus trust is the mechanism underlying
	\cite{ref16,ref17,ref18,ref19,ref20}: rather than a single Trusted
	Authority verifying every entity, a blockchain ledger distributes the
	verification record across Roadside Units or Trusted Authentication
	Servers, so that compromising one node does not compromise the
	authentication record for the network as a whole. Abbas et al.
	\cite{ref16} frame this as trading a single point of failure for a
	coordination overhead -- the recurring cost of removing centralized
	trust.
	
	\begin{figure}[t]
		\centering
		\fbox{\parbox[c][5.5cm][c]{8cm}{\centering\itshape
				Placeholder --- flow diagram of the five threat/trust mechanisms
				(injection/spoofing, eavesdropping/MITM, resource exhaustion,
				key-agreement authentication, ledger-based consensus) mapped
				against network layer, to be inserted.}}
		\caption{Threat and trust-establishment mechanisms mapped to
			vehicular network layer.}
		\label{fig:mechanisms}
	\end{figure}
	
	\subsection{Trust Establishment and Access-Control Techniques}
	
	Establishing entity trust in a resource-constrained, latency-sensitive
	vehicular network is achieved through three broad technique families,
	differing in where the trust anchor resides and what infrastructure
	they require: lightweight cryptographic techniques, decentralized
	ledger-based techniques, and hardware-rooted/selective techniques.
	
	Lightweight cryptographic techniques are the simplest and most widely
	deployable, requiring no additional network infrastructure beyond the
	communicating entities themselves. Dai et al. \cite{ref13} and Liu et
	al. \cite{ref14} fall in this category: security derives entirely from
	the mathematical hardness of the underlying problem (lattice-based
	KEMs, formally-verified handshake protocols) and validation is
	correspondingly fast -- \cite{ref13} reports roughly 60\% lower
	handshake latency than RSA. The limitation is that a compromised or
	lost key directly compromises the entity's identity, with no
	distributed fallback.
	
	Decentralized ledger-based techniques address that single-point-of-compromise
	limitation by distributing the authentication record itself. Rahayu et
	al. \cite{ref17}, Ali et al. \cite{ref18}, Kebande et al. \cite{ref19},
	and Li et al. \cite{ref20} all anchor authentication state in a
	blockchain accessible to multiple Roadside Units or Trusted
	Authentication Servers, so no single compromised node can forge or
	erase an authentication record. The cost is coordination and consensus
	latency: \cite{ref17} shows this latency scales with the number of
	Trusted Authentication Servers in a given topology, and \cite{ref18}
	reports the resulting computational overhead at under 2\% -- an
	acceptable but non-zero cost of removing centralization.
	
	Hardware-rooted and selective techniques target the specific
	weaknesses of the two families above. Madububa Mbachu et al.
	\cite{ref22} anchor trust in silicon-level manufacturing variation via
	a PUF rather than a storable secret, evaluated on real FPGA hardware
	with sub-40-microsecond encryption times. Luo et al. \cite{ref29} and
	Sheik et al. \cite{ref30} instead operate one layer up, applying
	selective, risk-prioritized authentication and formal
	threat-requirement matching rather than uniform enforcement everywhere
	-- directing scarce cryptographic and verification effort toward the
	highest-DREAD-scored threats rather than treating all links
	identically.
	
	\begin{table}[t]
		\centering
		\caption{Comparison of Vehicular Trust-Establishment Technique Families}
		\label{tab:trust_families}
		\small
		\begin{tabular}{p{1.7cm}p{1.6cm}p{1.6cm}p{2.4cm}}
			\toprule
			\textbf{Technique Family} & \textbf{Trust Anchor} & \textbf{Infrastructure} & \textbf{Key Trade-off} \\
			\midrule
			Lightweight cryptographic & Shared/derived session key & None beyond endpoints & Low latency vs.\ no distributed fallback \\
			\addlinespace
			Decentralized ledger-based & Distributed ledger record & RSU/TAS blockchain network & No single point of failure vs.\ consensus overhead \\
			\addlinespace
			Hardware-rooted / selective & Physical fingerprint or risk score & FPGA/PUF or formal threat model & Unclonability / focused effort vs.\ fabrication or modeling overhead \\
			\bottomrule
		\end{tabular}
	\end{table}
	
	As Table~\ref{tab:trust_families} summarizes, no single technique
	family dominates: lightweight cryptography optimizes for latency,
	ledger-based techniques optimize for elimination of centralized trust,
	and hardware-rooted/selective techniques optimize for physical
	unclonability and risk-proportional effort -- each accepting a
	different overhead in exchange. Representative references for each
	family are given in Table~\ref{tab:trust_families_refs}.
	
	\begin{table}[t]
		\centering
		\caption{Representative References by Technique Family}
		\label{tab:trust_families_refs}
		\small
		\begin{tabular}{p{3.6cm}p{3.6cm}}
			\toprule
			\textbf{Technique Family} & \textbf{Representative References} \\
			\midrule
			Lightweight cryptographic & \cite{ref13,ref14,ref15} \\
			Decentralized ledger-based & \cite{ref16,ref17,ref18,ref19,ref20} \\
			Hardware-rooted / selective & \cite{ref22,ref29,ref30} \\
			\bottomrule
		\end{tabular}
	\end{table}
	
	% ============================================================
	% SECTION III onward to be appended here as they are drafted.
	% ============================================================
	
	\begin{thebibliography}{99}
		
		\bibitem{ref1} M. Girdhar, Y. You, T.-J. Song, S. Ghosh, and J. Hong,
		``Post-Accident Cyberattack Event Analysis for Connected and Automated
		Vehicles,'' \emph{IEEE Access}, vol. 10, pp. 83176--XXXXX, 2022, doi:
		10.1109/ACCESS.2022.3196346.
		
		\bibitem{ref2} M. Girdhar, J. Hong, and J. Moore, ``Cybersecurity of
		Autonomous Vehicles: A Systematic Literature Review of Adversarial Attacks
		and Defense Models,'' \emph{IEEE Open J. Veh. Technol.}, vol. 4, 2023,
		doi: 10.1109/OJVT.2023.3265363.
		
		\bibitem{ref3} C. V. Kifor and A. Popescu, ``Automotive Cybersecurity: A
		Survey on Frameworks, Standards, and Testing and Monitoring
		Technologies,'' \emph{Sensors}, vol. 24, no. 18, p. 6139, 2024, doi:
		10.3390/s24186139.
		
		\bibitem{ref4} A. Abu Elkhail, R. U. D. Refat, R. Habre, A. Hafeez, A.
		Bacha, and H. Malik, ``Vehicle Security: A Survey of Security Issues and
		Vulnerabilities, Malware Attacks and Defenses,'' \emph{IEEE Access}, vol.
		9, pp. 162401--XXXXX, 2021, doi: 10.1109/ACCESS.2021.3130495.
		
		\bibitem{ref5} F. Luo, X. Zhang, Z. Yang, Y. Jiang, J. Wang, M. Wu, and W.
		Feng, ``Cybersecurity Testing for Automotive Domain: A Survey,''
		\emph{Sensors}, vol. 22, no. 23, p. 9211, 2022, doi: 10.3390/s22239211.
		
		\bibitem{ref6} O. Y. Al-Jarrah, C. Maple, M. Dianati, D. Oxtoby, and A.
		Mouzakitis, ``Intrusion Detection Systems for Intra-Vehicle Networks: A
		Review,'' \emph{IEEE Access}, vol. 7, pp. 21266--21289, 2019, doi:
		10.1109/ACCESS.2019.2894183.
		
		\bibitem{ref7} A. Ghadi and T. Khan Mohd, ``Learning-Based Intrusion
		Detection Systems for In-Vehicle CAN Networks: A Comprehensive Survey with
		Deployment and Real-Time Considerations,'' \emph{J. Braz. Comput. Soc.},
		vol. 32, p. 1, 2026, doi: 10.5753/jbcs.2026.6140.
		
		\bibitem{ref8} M. Almehdhar, A. Albaseer, M. A. Khan, M. Abdallah, H.
		Menouar, S. Al-Kuwari, and A. Al-Fuqaha, ``Deep Learning in the Fast Lane:
		A Survey on Advanced Intrusion Detection Systems for Intelligent Vehicle
		Networks,'' \emph{IEEE Open J. Veh. Technol.}, vol. 5, 2024, doi:
		10.1109/OJVT.2024.3422253.
		
		\bibitem{ref9} G. Karopoulos, G. Kambourakis, E. Chatzoglou, J. L.
		Hern\'andez-Ramos, and V. Kouliaridis, ``Demystifying In-Vehicle Intrusion
		Detection Systems: A Survey of Surveys and a Meta-Taxonomy,''
		\emph{Electronics}, vol. 11, no. 7, p. 1072, 2022, doi:
		10.3390/electronics11071072.
		
		\bibitem{ref10} M. Bozdal, M. Samie, S. Aslam, and I. Jennions,
		``Evaluation of CAN Bus Security Challenges,'' \emph{Sensors}, vol. 20,
		no. 8, p. 2364, 2020, doi: 10.3390/s20082364.
		
		\bibitem{ref11} T. C. M. D\"onmez, ``Anomaly Detection in Vehicular CAN Bus
		Using Message Identifier Sequences,'' \emph{IEEE Access}, vol. 9, pp.
		136243--XXXXX, 2021, doi: 10.1109/ACCESS.2021.3117038.
		
		\bibitem{ref12} A. Adu-Kyere, E. Nigussie, and J. Isoaho, ``Self-Aware
		Cybersecurity Architecture for Autonomous Vehicles: Security through
		System-Level Accountability,'' \emph{Sensors}, vol. 23, no. 21, p. 8817,
		2023, doi: 10.3390/s23218817.
		
		\bibitem{ref13} Y. Dai, Q. Wang, X. Song, and S. Wang, ``A Lightweight
		Key Agreement Protocol for V2X Communications Based on Kyber and
		Saber,'' \emph{Sensors}, vol. 25, no. 22, p. 6938, 2025, doi:
		10.3390/s25226938.
		
		\bibitem{ref14} Y. Liu, C. Li, Q. Sun, and H. Jiang, ``Security
		Authentication Scheme for Vehicle-to-Everything Computing Task
		Offloading Environments,'' \emph{Sensors}, vol. 25, no. 20, p. 6428,
		2025, doi: 10.3390/s25206428.
		
		\bibitem{ref15} P. Escapa Gord\'on, V. Matell\'an Olivera, and A.
		Su\'arez Corona, ``Vehicle-to-Vehicle Secure Communication Protocol
		Based on Identity-Based Cryptography Using the Vehicle Identification
		Number,'' \emph{Sensors}, vol. 25, no. 19, p. 5954, 2025, doi:
		10.3390/s25195954. \emph{[Title reconstructed from truncated OCR text
			-- verify exact wording on Scopus.]}
		
		\bibitem{ref16} S. Abbas, M. Abu Talib, A. Ahmed, F. Khan, S. Ahmad,
		and D.-H. Kim, ``Blockchain-Based Authentication in Internet of
		Vehicles: A Survey,'' \emph{Sensors}, 2024. \emph{[DOI/vol./issue
			pending Scopus verification.]}
		
		\bibitem{ref17} M. Rahayu, M. B. Hossain, S. Huda, Y. Kodera, M. A.
		Ali, and Y. Nogami, ``The Design and Implementation of
		Kerberos-Blockchain Vehicular Authentication for VANETs,''
		\emph{Sensors}, vol. 24, no. 23, p. 7428, 2024, doi:
		10.3390/s24237428.
		
		\bibitem{ref18} G. Ali, M. ElAffendi, and N. Ahmad, ``BlockAuth: A
		Blockchain-Based Framework for Secure Vehicle Authentication and
		Authorization,'' \emph{PLoS ONE}, vol. 18, no. 9, p. e0291596, 2023,
		doi: 10.1371/journal.pone.0291596.
		
		\bibitem{ref19} V. R. Kebande, F. M. Awaysheh, R. A. Ikuesan, S. A.
		Alawadi, and M. D. Alshehri, ``A Blockchain-Based Multi-Factor
		Authentication Model for Vehicular Cloud and Cloud-Enabled Internet of
		Vehicles.'' \emph{[Journal/vol./DOI pending Scopus verification.]}
		
		\bibitem{ref20} J. Li, Y. Lin, Y. Li, Y. Zhuang, and Y. Cao, ``BPA: A
		Novel Blockchain-Based Privacy-Preserving Authentication Scheme for
		the Internet of Vehicles,'' \emph{Electronics}, 2024. \emph{[Vol.
			/issue/DOI pending Scopus verification.]}
		
		\bibitem{ref21} Y. Wang, Y. Xu, Z. Liu, S. Liu, and Y. Wu, ``Research
		on Lightweight Dynamic Security Protocol for Intelligent In-Vehicle
		CAN Bus,'' \emph{Sensors}, vol. 25, no. 11, p. 3380, 2025, doi:
		10.3390/s25113380.
		
		\bibitem{ref22} U. Madububa Mbachu, R. Fatima, A. Sherif, E. Dockery,
		M. Mahmoud, M. Alsabaan, and K. Khalil, ``Hardware Acceleration-Based
		PUF Authentication for the Internet of Vehicles,'' \emph{Sensors},
		vol. 25, no. 5, p. 1629, 2025, doi: 10.3390/s25051629.
		
		\bibitem{ref23} E. F. Morgan and M. H. Ali, ``Digital Twin-Driven
		Cybersecurity for 5G/6G-Enabled Electric Vehicle Charging
		Infrastructure: A Review.'' \emph{[Journal/vol./DOI pending Scopus
			verification.]}
		
		\bibitem{ref24} S. Hamdare, O. Kaiwartya, M. Aljaidi, M. Jugran, Y.
		Cao, S. Kumar, M. Mahmud, D. Brown, and J. Lloret, ``Cybersecurity
		Risk Analysis of Electric Vehicle Charging Station Network
		Infrastructure.'' \emph{[Journal/vol./DOI pending Scopus
			verification.]}
		
		\bibitem{ref25} J. Johnson, T. Berg, B. Anderson, and B. Wright,
		``Review of Electric Vehicle Charger Cybersecurity Vulnerabilities,
		Potential Impacts, and Defenses,'' \emph{Sensors}, vol. 23, no. 15,
		p. 6716, 2023, doi: 10.3390/s23156716.
		
		\bibitem{ref26} S. Hamdare, D. J. Brown, Y. Cao, M. Aljaidi, S.
		Kumar, R. Alanazi, M. Jugran, P. Vyas, and O. Kaiwartya, ``A Novel
		Charging Management and Security Framework for the Electric Vehicle
		Ecosystem.'' \emph{[Journal/vol./DOI pending Scopus verification.]}
		
		\bibitem{ref27} M. ElKashlan, M. S. Elsayed, A. D. Jurcut, and M.
		Azer, ``A Machine Learning-Based Intrusion Detection System for IoT
		Electric Vehicle Charging Stations.'' \emph{[Journal/vol./DOI
			pending Scopus verification.]}
		
		\bibitem{ref28} Y. Wang, Y. Wu, Y. Xu, K. Zhang, and Y. Xu, ``Research
		on Network Intrusion Detection Based on Weighted Histogram Algorithm
		for In-Vehicle Ethernet.'' \emph{[Journal/vol./DOI pending Scopus
			verification.]}
		
		\bibitem{ref29} F. Luo, Y. Jiang, J. Wang, Z. Li, and X. Zhang, ``A
		Framework for Cybersecurity Requirements Management in the Automotive
		Domain.'' \emph{[Journal/vol./DOI pending Scopus verification.]}
		
		\bibitem{ref30} A. T. Sheik, C. Maple, G. Epiphaniou, and M. Dianati,
		``Securing Cloud-Assisted Connected and Autonomous Vehicles: An
		In-Depth Threat Analysis and Risk Assessment.'' \emph{[Journal/vol./DOI
			pending Scopus verification.]}
		
		\bibitem{refDini} G. Dini \emph{et al.}, [Embedded ML-based voltage
		fingerprinting for CAN bus ECU identification], \emph{IEEE Access}, 2025,
		doi: 10.1109/ACCESS.2025.3545245.
		
	\end{thebibliography}
	
\end{document}